\chapter{Data Acquisition, Preparation and Feature Engineering}
\label{chap:data}

\section{Source, provenance and evidence boundary}

Sufficient representative plant history was not available for this initiation project, so the final workflow uses one supplied deterministic synthetic Comma-Separated Values (CSV) file. It spans 16 April--16 July 2026 and contains exactly 1,589,760 unique timestamps on a gap-free five-second grid. The three laboratory fields are populated at 1,104 timestamps, equivalent to a regular two-hour pattern. Notebook 01 audits the source; Notebook 02 produces the causal supervised handoff; Notebooks 03 and 04 export the final model artifacts.

The word \emph{deterministic} needs a precise boundary. The file hash matches its manifest, which makes the supplied file reproducible as an artifact. However, the repository does not contain the original generator code, configuration path or random seed. I could therefore audit the data actually present, but I could not reproduce or verify exact generator equations. Relationships reported in this chapter are empirical properties of the CSV, not recovered plant laws or undocumented generation formulas.

\begin{table}[ht]
\centering\small
\caption{Audited data products and their distinct roles}
\label{tab:datasets}
\begin{tabularx}{\textwidth}{@{}p{0.25\textwidth}rrrrY@{}}
\toprule
Data product & Rows & Columns & Lab rows & Cadence & Role \\
\midrule
Canonical source & 1,589,760 & 17 & 1,104 & 5 s / 2 h lab & Common source for audit and model development. \\
Aligned supervised handoff & 1,103 & 18 & 1,103 & 2 h & Timestamp, 16 inputs and final-moisture target. \\
Canonical TEST interval & 237,600 & 18 & 165 & 5 s / 2 h lab & Frozen-model analytical and database evaluation. \\
Dashboard demo fork & 34,560 & 18 & preserved & 5 s / 2 h lab & Two-day jury replay with disclosed injected episodes; not model evaluation. \\
\bottomrule
\end{tabularx}
\source{Direct repository audit generated for this report.}
\end{table}

The source labels form 22 synthetic disturbance episodes: five steam dips, five airflow restrictions, four cooling losses, four fan-speed dips and four feed surges. These labels are used for the episode counts and held-out analysis.

\begin{table}[ht]
\centering\small
\caption{Disturbance episodes found directly in the canonical CSV}
\label{tab:scenario-counts}
\begin{tabular}{@{}lrr@{}}
\toprule
Label & Episodes & Interpretation \\
\midrule
Steam dip & 5 & Synthetic reduction in thermal-supply proxy. \\
Airflow restriction & 5 & Synthetic drying-air disturbance. \\
Cooling loss & 4 & Synthetic downstream cooling disturbance. \\
Fan-speed dip & 4 & Synthetic air/exhaust equipment disturbance. \\
Feed surge & 4 & Synthetic dryer-loading disturbance. \\
\midrule
Total & 22 & Label audit; not confirmed plant failures. \\
\bottomrule
\end{tabular}
\end{table}

The prototype uses nine process inputs and three sparse laboratory fields. Missing laboratory timestamps remain blank so that an old result is never presented as a fresh measurement. Vacuum is retained using the convention of the supplied project data; plant tag polarity should be confirmed before industrial deployment \cite{datadictionary}.

\section{Exploration without over-interpreting the generator}

The process variables occupy compact synthetic ranges with gradual variation and labelled disturbances. Direct calculation gives a fan-speed/air-flow correlation of 0.938, an air-flow/vacuum correlation of 0.979 and a feed/residence correlation of $-0.935$. At laboratory timestamps, dryer temperature and final moisture correlate at $-0.508$, while cooler temperature and final product temperature correlate at 0.909. These patterns are useful for checking consistency, but they should not be read as causal plant relationships.

I kept unusual values unless a specific data-quality rule rejected them. In a supervision project, a rare point may be a valid transient, a designed disturbance or an acquisition error. Representative historian data would also need operating mode and sensor-quality information that the synthetic file does not contain.

\begin{figure}[p]
\centering
\includegraphics[width=.95\textwidth]{data_distributions.png}
\caption{Distributions of process and sparse quality variables in the canonical source.}
\source{Executed Notebook 01 output.}
\label{fig:eda-distributions}
\end{figure}

\begin{figure}[p]
\centering
\includegraphics[width=.94\textwidth]{correlation_matrix.png}
\caption{Pearson correlation matrix for canonical process and sparse quality variables.}
\source{Executed Notebook 02 output.}
\label{fig:eda-correlation}
\end{figure}

\FloatBarrier

\section{Causal process-to-laboratory alignment}

The main leakage risk is not a random missing value; it is time. A laboratory result at $t_{lab}$ describes material that passed through the dryer earlier. The direct process snapshot is selected at
\begin{equation}
t_{effective}=t_{lab}-t_{residence}(t_{lab})-t_{transport}.
\end{equation}
The generated artifact uses the most recent process row at or before $t_{effective}$ and a configurable transport delay of zero minutes. Density and final product temperature come from the most recent laboratory row strictly before $t_{lab}$. Final moisture is the target and never an input.

\begin{lstlisting}[language=Python,caption={Causal snapshot and previous-laboratory selection.}]
effective_time = effective_window_end(
    lab_timestamp, residence, transport_delay_minutes
)
process_row = process.loc[
    process["Timestamp"] <= effective_time
].iloc[-1]
previous_lab = laboratory.loc[
    laboratory["Sample Timestamp"] < lab_timestamp
].iloc[-1]
\end{lstlisting}

The first laboratory row has no strictly previous laboratory context, leaving 1,103 supervised observations. The chronology is split into 772 TRAIN, 166 VALIDATION and 165 TEST rows. Five expanding folds tune models inside TRAIN, the later VALIDATION block selects the family, and TEST remains closed until selection is complete. Blocked evaluation is appropriate here because random splitting can mix temporally dependent process regimes and make future information indirectly available \cite{bergmeir2012}.

\section{Feature engineering and dimensional meaning}

The final moisture contract contains nine direct process values, two prior laboratory values and five interactions. The anomaly contract contains the nine process values plus six instant interactions. Standardization is learned from training data only. The main design choices are summarized below.

\begin{table}[htbp]
\centering\small
\caption{Retained engineered features and what they do not represent}
\label{tab:engineered-summary}
\begin{tabularx}{\textwidth}{@{}p{0.27\textwidth}p{0.27\textwidth}Y@{}}
\toprule
Feature family & Construction & Engineering boundary \\
\midrule
Temperature differences & Dryer air minus inlet, cooler air or prior final-product temperature & Thermal contrast; not heat flow. \\
Air per feed & Air flow divided by wet-product feed & Relative gas/loading index; dimensionless only under the shared volumetric convention. \\
Steam per feed & Steam pressure divided by wet-product feed & Statistical supply/loading index; units bar/(m$^3$/h). \\
Steam--temperature interaction / thermal load & Steam pressure multiplied by dryer-air temperature & Regression interaction; not thermal duty. \\
Heating index / thermal exposure & Residence time multiplied by dryer-air temperature & Time--temperature exposure proxy; not energy. \\
\bottomrule
\end{tabularx}
\end{table}

Feature names are deliberately qualified because a plausible label can otherwise overstate the underlying physics. The present data do not support a closed mass/energy balance. At runtime, protected denominators and required-input checks produce a \code{DATA\_QUALITY} or unavailable state rather than a silent infinity or improvised prediction. The saved schemas lock feature names, order and versions for both models.

\begin{engineeringbox}[Data conclusion]
The canonical source is adequate for reproducible software and method testing. Its size does not remove the central limitation: only 1,103 rows contain supervised moisture targets, the disturbance labels are synthetic, and the original generation procedure is unavailable.
\end{engineeringbox}
